% !TeX program = xelatex
\documentclass[12pt,a4paper]{article}
\usepackage{fontspec}
\setmainfont{Times New Roman}
\usepackage[vietnamese]{babel}
\usepackage{amsmath,amssymb}
\usepackage[margin=2.3cm]{geometry}
\usepackage{enumitem}
\setlist[enumerate]{leftmargin=*,itemsep=0.6em}
\setlength{\parindent}{0pt}
\begin{document}
\begin{center}
{\Large\bfseries ĐÁP ÁN VÀ LỜI GIẢI -- ĐỀ 2 (2025)}\\[2pt]
{\large\bfseries PHẦN ĐẠI SỐ VÀ TOÁN ỨNG DỤNG}
\end{center}

\section*{Câu I. Thống kê và xác suất}
\begin{enumerate}
\item Nhóm $[12;16)$ có tần số là 75. Tần số tương đối của nhóm là
\[
\frac{75}{300}\cdot100\%=25\%.
\]
\item Không gian mẫu gồm 8 kết quả đồng khả năng: rút được thẻ mang một trong các số từ 1 đến 8. Các kết quả thuận lợi cho $A$ là rút được thẻ số 3 hoặc số 6, có 2 kết quả. Do đó
\[
P(A)=\frac{2}{8}=\frac14.
\]
\end{enumerate}

\section*{Câu II. Biểu thức đại số}
\begin{enumerate}
\item Với $x=9$ thì $\sqrt{x}=3$, nên
\[
A=\frac{3+2}{3-2}=5.
\]
\item Đặt $t=\sqrt{x}>0$, với $t\ne2$. Quy đồng hai phân thức trong $B$:
\begin{align*}
B&=\frac{t^2+t-4}{t(t-2)}-\frac{t}{t(t-2)}\\
&=\frac{t^2-4}{t(t-2)}
=\frac{(t-2)(t+2)}{t(t-2)}
=\frac{t+2}{t}
=\frac{\sqrt{x}+2}{\sqrt{x}}.
\end{align*}
Các phép biến đổi hợp lệ do $x>0$ và $x\ne4$.
\item Dùng kết quả trên,
\[
\frac AB=\frac{\frac{\sqrt{x}+2}{\sqrt{x}-2}}{\frac{\sqrt{x}+2}{\sqrt{x}}}=\frac{\sqrt{x}}{\sqrt{x}-2}.
\]
Do $\sqrt{x}+2>0$, bất đẳng thức tương đương
\[
\frac{\sqrt{x}}{\sqrt{x}-2}<\frac12
\iff \frac{2\sqrt{x}-(\sqrt{x}-2)}{2(\sqrt{x}-2)}<0
\iff \frac{\sqrt{x}+2}{2(\sqrt{x}-2)}<0.
\]
Suy ra $\sqrt{x}-2<0$, tức $x<4$. Kết hợp $x$ là số nguyên dương (và điều kiện xác định), các giá trị là $1,2,3$. Giá trị lớn nhất là $\boxed{3}$.
\end{enumerate}

\section*{Câu III. Bài toán thực tế và hệ thức Viète}
\begin{enumerate}
\item Gọi quãng đường Hà Nội--Hải Phòng là $d$ km ($d>0$). Thời gian đi là $d/60$ giờ, thời gian về là $d/40$ giờ. Theo đề bài,
\[
\frac d{40}-\frac d{60}=1\quad\Longrightarrow\quad \frac d{120}=1\quad\Longrightarrow\quad d=120.
\]
Vậy quãng đường dài $\boxed{120\text{ km}}$.
\item Gọi giá niêm yết của ba lô và máy tính cầm tay lần lượt là $a,b$ nghìn đồng ($a,b>0$). Ta có
\[
\begin{cases}
a+b=885,\\
0{,}8a+0{,}75b=682.
\end{cases}
\]
Từ phương trình đầu, $a=885-b$. Thay vào phương trình sau:
\[
0{,}8(885-b)+0{,}75b=682
\Longrightarrow 708-0{,}05b=682
\Longrightarrow b=520.
\]
Do đó $a=885-520=365$. Giá niêm yết lần lượt là $\boxed{365}$ nghìn đồng và $\boxed{520}$ nghìn đồng.
\item Theo Viète,
\[
x_1+x_2=-8,\qquad x_1x_2=-6.
\]
Vì $x_2$ là nghiệm khác 0 nên nhân điều kiện đã cho với $x_2$ được
\[
70-mx_1^2=x_1x_2+mx_2^2
\Longleftrightarrow m(x_1^2+x_2^2)+x_1x_2=70.
\]
Mà
\[
x_1^2+x_2^2=(x_1+x_2)^2-2x_1x_2=64+12=76.
\]
Vậy $76m-6=70$, suy ra $\boxed{m=1}$.
\end{enumerate}

\section*{Câu IV. Diện tích và thể tích}
\begin{enumerate}
\item Diện tích xung quanh hình trụ là
\[
S_{\text{xq}}=2\pi rh\approx2\cdot3{,}14\cdot50\cdot150=47\,100\text{ cm}^2.
\]
\item Phần nước đã dùng có chiều cao 40 cm, nên thể tích là
\[
V=\pi r^2h\approx3{,}14\cdot50^2\cdot40=314\,000\text{ cm}^3.
\]
\end{enumerate}

\section*{Câu V. Bài toán tối ưu}
Gọi $n$ là số xe được bổ sung; $n$ là số nguyên không âm. Khi đó lợi nhuận trung bình mỗi xe mỗi ngày là $1-0{,}02n$ triệu đồng, tổng số xe là $35+n$, và tổng lợi nhuận là
\[
M=(35+n)(1-0{,}02n)=-0{,}02n^2+0{,}3n+35.
\]
Hoàn thành bình phương:
\[
M=-0{,}02(n-7{,}5)^2+36{,}125.
\]
Với $n$ nguyên không âm, biểu thức đạt lớn nhất khi $n$ gần $7{,}5$ nhất, tức $n=7$ hoặc $n=8$. Hai phương án đều cho $M=36{,}12$ triệu đồng/ngày. Công ty nên bổ sung $\boxed{7\text{ hoặc }8\text{ xe}}$.
\end{document}
